% Mechanistic Power Check and Occupation-Level Aggregation After Authoritative Relabeling
% This file validates that task-level heterogeneity is sufficient for identifiability

After upgrading $(s_j,c_j)$ from subjective keyword tagging to an "authoritative evaluations $\rightarrow$ capability dimensions $\rightarrow$ semantic matching $\rightarrow$ piecewise assignment" pipeline,
we must validate \textbf{mechanistic power}: task-level heterogeneity in $(s_j,c_j)$ should be sufficient, and task weights $w_j$ should not collapse $(\alpha,\beta)$ into a single default-driven outcome.
Otherwise, even authoritative sources would yield weak identifiability in the evolution model.

\subsection{Occupation-Level Shock Aggregates}

For the usable task set (size $N$, allowing $N<20$) renormalized such that $\sum_{j=1}^{N} w_j=1$, we define
\begin{equation}
    \alpha=\sum_{j=1}^{N} w_j s_j,\qquad
    \beta=\sum_{j=1}^{N} w_j c_j,\qquad
    \mu=\beta-\alpha.
\end{equation}

\subsection{Diagnostics: Heterogeneity and Effective Diversity}

\paragraph{(1) Weighted variance (heterogeneity)}
\begin{equation}
    \mathrm{Var}_w(s)=\sum_{j=1}^{N} w_j (s_j-\alpha)^2,\qquad
    \mathrm{Var}_w(c)=\sum_{j=1}^{N} w_j (c_j-\beta)^2.
\end{equation}
Low variance indicates compressed scoring or overly convergent mapping rules, weakening identifiability.

\paragraph{(2) Weighted tier histograms (avoid single-tier dominance)}
Let $\mathcal{G}_{LMH}=\{0.25,0.5,0.75\}$. Define tier weight mass
\begin{equation}
    h_s(g)=\sum_{j:\,s_j=g} w_j,\qquad
    h_c(g)=\sum_{j:\,c_j=g} w_j,\qquad g\in\mathcal{G}_{LMH}.
\end{equation}
Large $\max_g h_s(g)$ or $\max_g h_c(g)$ flags default-driven collapse.

\paragraph{(3) Weight mass by capability dimension (avoid weight concentration)}
For each task $j$, let $\mathrm{dim}_j$ be its mapped dimension set (1--2 dimensions).
If multi-mapped, let $\pi_{j d}$ allocate task $j$ to dimension $d$ with $\sum_{d\in\mathrm{dim}_j}\pi_{j d}=1$.
Dimension weight mass is
\begin{equation}
    m(d)=\sum_{j=1}^{N} w_j\,\pi_{j d}.
\end{equation}
High $\max_d m(d)$ indicates low effective diversity despite broad nominal coverage.

\subsection{Engineering Thresholds and Actions (Gatekeeping Before $E(t)$)}

We use engineering thresholds as a gate to decide whether to proceed to scenario evolution:
\begin{equation}
\mathrm{Var}_w(s)\ge \varepsilon_s,\quad
\mathrm{Var}_w(c)\ge \varepsilon_c,\quad
\max_g h_s(g)\le \eta_s,\quad
\max_d m(d)\le \eta_d,
\end{equation}
with defaults $\varepsilon_s=\varepsilon_c=0.005$, $\eta_s=0.70$, and $\eta_d=0.60$.
Violations imply high mechanism risk, requiring refinement of tiering or the dimension-to-$(s,c)$ mapping (without changing $w_j$) before advancing to scenario-driven $E(t)$ dynamics.
